% Persistas & Partners — firm letterhead.
%
% This is the firm's identity as it appears on every piece of correspondence
% that leaves the office: the mark, the two partners with their enrolment
% details, and the registered office in the footer. It is separated from
% persist-base.tex because a template may want the firm's typography without
% its letterhead — an internal memo, a draft for comment.
%
% % Persist — shared document preamble.
%
% Every template begins with % Persist — shared document preamble.
%
% Every template begins with % Persist — shared document preamble.
%
% Every template begins with \input{_shared/persist-base}. The PRD lists
% nineteen templates for the Smart Form Compiler; each carrying its own copy of
% the font, colour and geometry setup would mean nineteen places to change when
% the firm's letterhead does, and nineteen chances to miss one.
%
% ENGINE: XeLaTeX. Not pdfLaTeX — see src-tauri/KEEL-RULES.md.
%
% Templates get: A4 geometry, Noto (with Devanagari declared), the Persist
% colour tokens, the L/C/R column types, and \firmfooter. They add only their
% own content.

\usepackage[top=20mm, bottom=20mm, left=22mm, right=22mm]{geometry}

% ── Fonts ────────────────────────────────────────────────────────────────────
% Noto covers ₹ (U+20B9) and Devanagari. pdfLaTeX fails outright on the former.
\usepackage{fontspec}
\setmainfont{Noto Serif}[
  UprightFont = *,
  BoldFont    = * Bold,
  ItalicFont  = * Italic,
]
\setsansfont{Noto Sans}
\setmonofont{Noto Sans Mono}[Scale=MatchLowercase]
% Declared so a Hindi party name or title renders rather than dropping to tofu.
% Devanagari needs its own shaper, not merely a font that has the glyphs.
\newfontfamily\devanagarifont{Noto Serif Devanagari}[Script=Devanagari]
% Wrap Devanagari runs: {\hindi ...}
\newcommand{\hindi}[1]{{\devanagarifont #1}}

% ── Packages ─────────────────────────────────────────────────────────────────
\usepackage{microtype}
\usepackage{booktabs}
\usepackage{tabularx}
\usepackage{array}
\usepackage{longtable}
\usepackage{colortbl}
\usepackage{xcolor}
\usepackage{graphicx}
\usepackage{fancyhdr}
\usepackage{multirow}
\usepackage{hhline}
\usepackage{calc}
\usepackage{enumitem}
\usepackage[hidelinks]{hyperref}
\usepackage{parskip}

% ── Colours — mirror src/design-system/tokens.ts ─────────────────────────────
% A client sees the invoice and the portal side by side; they should look like
% the same firm.
\definecolor{accentBlue}{RGB}{74,101,128}     % --color-accent-primary
\definecolor{terracotta}{RGB}{181,96,74}      % --color-accent-secondary
\definecolor{textPrimary}{RGB}{44,44,42}      % --color-text-primary
\definecolor{textSecondary}{RGB}{107,104,98}  % --color-text-secondary
\definecolor{textTertiary}{RGB}{154,149,144}  % --color-text-tertiary
\definecolor{borderGray}{RGB}{226,221,216}    % --color-border
\definecolor{bgSecondary}{RGB}{240,236,229}   % --color-bg-secondary
\definecolor{statusGreen}{RGB}{74,124,89}     % --color-status-clear
\definecolor{statusUrgent}{RGB}{192,57,43}    % --color-status-urgent

% ── Typography ───────────────────────────────────────────────────────────────
\renewcommand{\familydefault}{\rmdefault}
\setlength{\parindent}{0pt}

% ── Fixed-width column types ─────────────────────────────────────────────────
\newcolumntype{R}[1]{>{\raggedleft\arraybackslash}p{#1}}
\newcolumntype{L}[1]{>{\raggedright\arraybackslash}p{#1}}
\newcolumntype{C}[1]{>{\centering\arraybackslash}p{#1}}

% ── Firm footer ──────────────────────────────────────────────────────────────
% Takes the firm identity as arguments rather than reading {{PLACEHOLDERS}}:
% a shared file that assumed particular placeholder names would force every
% template to declare them, whether or not the document shows them.
%   \firmfooter{name}{gstin}{pan}
\newcommand{\firmfooter}[3]{%
  \pagestyle{fancy}%
  \fancyhf{}%
  \renewcommand{\headrulewidth}{0pt}%
  \fancyfoot[C]{\footnotesize\color{textSecondary}
    Page \thepage\ — #1 — GSTIN: #2 — PAN: #3}%
}

% A footer for documents that are not invoices: no tax identifiers, just the
% firm and the page.
%   \plainfooter{name}
\newcommand{\plainfooter}[1]{%
  \pagestyle{fancy}%
  \fancyhf{}%
  \renewcommand{\headrulewidth}{0pt}%
  \fancyfoot[C]{\footnotesize\color{textSecondary} Page \thepage\ — #1}%
}

% ── Letterhead ───────────────────────────────────────────────────────────────
%   \letterhead{firm name}{address}{contact line}
\newcommand{\letterhead}[3]{%
  {\LARGE\bfseries\color{accentBlue} #1}\par
  \vspace{2pt}
  {\small\color{textSecondary} #2}\par
  {\small\color{textSecondary} #3}\par
  \vspace{2mm}
  {\color{borderGray}\rule{\textwidth}{0.4pt}}
  \vspace{4mm}
}

% ── Section heading in the firm's style ──────────────────────────────────────
\newcommand{\persistsection}[1]{%
  \vspace{3mm}
  {\small\bfseries\color{textSecondary} \MakeUppercase{#1}}\par
  \vspace{1mm}
}
. The PRD lists
% nineteen templates for the Smart Form Compiler; each carrying its own copy of
% the font, colour and geometry setup would mean nineteen places to change when
% the firm's letterhead does, and nineteen chances to miss one.
%
% ENGINE: XeLaTeX. Not pdfLaTeX — see src-tauri/KEEL-RULES.md.
%
% Templates get: A4 geometry, Noto (with Devanagari declared), the Persist
% colour tokens, the L/C/R column types, and \firmfooter. They add only their
% own content.

\usepackage[top=20mm, bottom=20mm, left=22mm, right=22mm]{geometry}

% ── Fonts ────────────────────────────────────────────────────────────────────
% Noto covers ₹ (U+20B9) and Devanagari. pdfLaTeX fails outright on the former.
\usepackage{fontspec}
\setmainfont{Noto Serif}[
  UprightFont = *,
  BoldFont    = * Bold,
  ItalicFont  = * Italic,
]
\setsansfont{Noto Sans}
\setmonofont{Noto Sans Mono}[Scale=MatchLowercase]
% Declared so a Hindi party name or title renders rather than dropping to tofu.
% Devanagari needs its own shaper, not merely a font that has the glyphs.
\newfontfamily\devanagarifont{Noto Serif Devanagari}[Script=Devanagari]
% Wrap Devanagari runs: {\hindi ...}
\newcommand{\hindi}[1]{{\devanagarifont #1}}

% ── Packages ─────────────────────────────────────────────────────────────────
\usepackage{microtype}
\usepackage{booktabs}
\usepackage{tabularx}
\usepackage{array}
\usepackage{longtable}
\usepackage{colortbl}
\usepackage{xcolor}
\usepackage{graphicx}
\usepackage{fancyhdr}
\usepackage{multirow}
\usepackage{hhline}
\usepackage{calc}
\usepackage{enumitem}
\usepackage[hidelinks]{hyperref}
\usepackage{parskip}

% ── Colours — mirror src/design-system/tokens.ts ─────────────────────────────
% A client sees the invoice and the portal side by side; they should look like
% the same firm.
\definecolor{accentBlue}{RGB}{74,101,128}     % --color-accent-primary
\definecolor{terracotta}{RGB}{181,96,74}      % --color-accent-secondary
\definecolor{textPrimary}{RGB}{44,44,42}      % --color-text-primary
\definecolor{textSecondary}{RGB}{107,104,98}  % --color-text-secondary
\definecolor{textTertiary}{RGB}{154,149,144}  % --color-text-tertiary
\definecolor{borderGray}{RGB}{226,221,216}    % --color-border
\definecolor{bgSecondary}{RGB}{240,236,229}   % --color-bg-secondary
\definecolor{statusGreen}{RGB}{74,124,89}     % --color-status-clear
\definecolor{statusUrgent}{RGB}{192,57,43}    % --color-status-urgent

% ── Typography ───────────────────────────────────────────────────────────────
\renewcommand{\familydefault}{\rmdefault}
\setlength{\parindent}{0pt}

% ── Fixed-width column types ─────────────────────────────────────────────────
\newcolumntype{R}[1]{>{\raggedleft\arraybackslash}p{#1}}
\newcolumntype{L}[1]{>{\raggedright\arraybackslash}p{#1}}
\newcolumntype{C}[1]{>{\centering\arraybackslash}p{#1}}

% ── Firm footer ──────────────────────────────────────────────────────────────
% Takes the firm identity as arguments rather than reading {{PLACEHOLDERS}}:
% a shared file that assumed particular placeholder names would force every
% template to declare them, whether or not the document shows them.
%   \firmfooter{name}{gstin}{pan}
\newcommand{\firmfooter}[3]{%
  \pagestyle{fancy}%
  \fancyhf{}%
  \renewcommand{\headrulewidth}{0pt}%
  \fancyfoot[C]{\footnotesize\color{textSecondary}
    Page \thepage\ — #1 — GSTIN: #2 — PAN: #3}%
}

% A footer for documents that are not invoices: no tax identifiers, just the
% firm and the page.
%   \plainfooter{name}
\newcommand{\plainfooter}[1]{%
  \pagestyle{fancy}%
  \fancyhf{}%
  \renewcommand{\headrulewidth}{0pt}%
  \fancyfoot[C]{\footnotesize\color{textSecondary} Page \thepage\ — #1}%
}

% ── Letterhead ───────────────────────────────────────────────────────────────
%   \letterhead{firm name}{address}{contact line}
\newcommand{\letterhead}[3]{%
  {\LARGE\bfseries\color{accentBlue} #1}\par
  \vspace{2pt}
  {\small\color{textSecondary} #2}\par
  {\small\color{textSecondary} #3}\par
  \vspace{2mm}
  {\color{borderGray}\rule{\textwidth}{0.4pt}}
  \vspace{4mm}
}

% ── Section heading in the firm's style ──────────────────────────────────────
\newcommand{\persistsection}[1]{%
  \vspace{3mm}
  {\small\bfseries\color{textSecondary} \MakeUppercase{#1}}\par
  \vspace{1mm}
}
. The PRD lists
% nineteen templates for the Smart Form Compiler; each carrying its own copy of
% the font, colour and geometry setup would mean nineteen places to change when
% the firm's letterhead does, and nineteen chances to miss one.
%
% ENGINE: XeLaTeX. Not pdfLaTeX — see src-tauri/KEEL-RULES.md.
%
% Templates get: A4 geometry, Noto (with Devanagari declared), the Persist
% colour tokens, the L/C/R column types, and \firmfooter. They add only their
% own content.

\usepackage[top=20mm, bottom=20mm, left=22mm, right=22mm]{geometry}

% ── Fonts ────────────────────────────────────────────────────────────────────
% Noto covers ₹ (U+20B9) and Devanagari. pdfLaTeX fails outright on the former.
\usepackage{fontspec}
\setmainfont{Noto Serif}[
  UprightFont = *,
  BoldFont    = * Bold,
  ItalicFont  = * Italic,
]
\setsansfont{Noto Sans}
\setmonofont{Noto Sans Mono}[Scale=MatchLowercase]
% Declared so a Hindi party name or title renders rather than dropping to tofu.
% Devanagari needs its own shaper, not merely a font that has the glyphs.
\newfontfamily\devanagarifont{Noto Serif Devanagari}[Script=Devanagari]
% Wrap Devanagari runs: {\hindi ...}
\newcommand{\hindi}[1]{{\devanagarifont #1}}

% ── Packages ─────────────────────────────────────────────────────────────────
\usepackage{microtype}
\usepackage{booktabs}
\usepackage{tabularx}
\usepackage{array}
\usepackage{longtable}
\usepackage{colortbl}
\usepackage{xcolor}
\usepackage{graphicx}
\usepackage{fancyhdr}
\usepackage{multirow}
\usepackage{hhline}
\usepackage{calc}
\usepackage{enumitem}
\usepackage[hidelinks]{hyperref}
\usepackage{parskip}

% ── Colours — mirror src/design-system/tokens.ts ─────────────────────────────
% A client sees the invoice and the portal side by side; they should look like
% the same firm.
\definecolor{accentBlue}{RGB}{74,101,128}     % --color-accent-primary
\definecolor{terracotta}{RGB}{181,96,74}      % --color-accent-secondary
\definecolor{textPrimary}{RGB}{44,44,42}      % --color-text-primary
\definecolor{textSecondary}{RGB}{107,104,98}  % --color-text-secondary
\definecolor{textTertiary}{RGB}{154,149,144}  % --color-text-tertiary
\definecolor{borderGray}{RGB}{226,221,216}    % --color-border
\definecolor{bgSecondary}{RGB}{240,236,229}   % --color-bg-secondary
\definecolor{statusGreen}{RGB}{74,124,89}     % --color-status-clear
\definecolor{statusUrgent}{RGB}{192,57,43}    % --color-status-urgent

% ── Typography ───────────────────────────────────────────────────────────────
\renewcommand{\familydefault}{\rmdefault}
\setlength{\parindent}{0pt}

% ── Fixed-width column types ─────────────────────────────────────────────────
\newcolumntype{R}[1]{>{\raggedleft\arraybackslash}p{#1}}
\newcolumntype{L}[1]{>{\raggedright\arraybackslash}p{#1}}
\newcolumntype{C}[1]{>{\centering\arraybackslash}p{#1}}

% ── Firm footer ──────────────────────────────────────────────────────────────
% Takes the firm identity as arguments rather than reading {{PLACEHOLDERS}}:
% a shared file that assumed particular placeholder names would force every
% template to declare them, whether or not the document shows them.
%   \firmfooter{name}{gstin}{pan}
\newcommand{\firmfooter}[3]{%
  \pagestyle{fancy}%
  \fancyhf{}%
  \renewcommand{\headrulewidth}{0pt}%
  \fancyfoot[C]{\footnotesize\color{textSecondary}
    Page \thepage\ — #1 — GSTIN: #2 — PAN: #3}%
}

% A footer for documents that are not invoices: no tax identifiers, just the
% firm and the page.
%   \plainfooter{name}
\newcommand{\plainfooter}[1]{%
  \pagestyle{fancy}%
  \fancyhf{}%
  \renewcommand{\headrulewidth}{0pt}%
  \fancyfoot[C]{\footnotesize\color{textSecondary} Page \thepage\ — #1}%
}

% ── Letterhead ───────────────────────────────────────────────────────────────
%   \letterhead{firm name}{address}{contact line}
\newcommand{\letterhead}[3]{%
  {\LARGE\bfseries\color{accentBlue} #1}\par
  \vspace{2pt}
  {\small\color{textSecondary} #2}\par
  {\small\color{textSecondary} #3}\par
  \vspace{2mm}
  {\color{borderGray}\rule{\textwidth}{0.4pt}}
  \vspace{4mm}
}

% ── Section heading in the firm's style ──────────────────────────────────────
\newcommand{\persistsection}[1]{%
  \vspace{3mm}
  {\small\bfseries\color{textSecondary} \MakeUppercase{#1}}\par
  \vspace{1mm}
}
 first; this builds on its colours and fonts.
%
% The header and footer are set once here rather than per template. When the
% firm moves office or a partner's number changes, this file changes and every
% document produced afterwards is correct.

\usepackage{lastpage}     % "Page 3 of 19" — the count, not just the number
\usepackage{tabularx}
\usepackage{changepage}   % adjustwidth, for indented section bodies
\usepackage{ragged2e}     % justifying

% ── Partner block ────────────────────────────────────────────────────────────
% One partner's line in the header: name, designation, phone, email.
%   \partnerblock{name}{designation}{phone}{email}
\newcommand{\partnerblock}[4]{%
  \begin{minipage}[t]{\linewidth}
    \raggedleft
    {\sffamily\bfseries\fontsize{7.5}{9}\selectfont #1}\\[1pt]
    {\sffamily\fontsize{7}{8.5}\selectfont #2}\\[1pt]
    {\sffamily\fontsize{7}{8.5}\selectfont #3}\\[1pt]
    {\sffamily\fontsize{7}{8.5}\selectfont #4}
  \end{minipage}%
}

% ── The letterhead itself ────────────────────────────────────────────────────
% Applied to every page, because a notice is served page by page and a page
% without the firm's mark is a page whose provenance is arguable.
%
% A template fills these in with \renewcommand and then calls \applyletterhead.
% Passing thirteen arguments to one macro would be unreadable at the call site
% and unreachable from fancyhdr, which rebuilds the bands on every page.
\newcommand{\firmlogo}{_shared/assets/persistas-logo.png}
\newcommand{\partneronename}{}   \newcommand{\partneronerole}{}
\newcommand{\partneronephone}{}  \newcommand{\partneroneemail}{}
\newcommand{\partnertwoname}{}   \newcommand{\partnertworole}{}
\newcommand{\partnertwophone}{}  \newcommand{\partnertwoemail}{}
\newcommand{\firmwebsite}{}      \newcommand{\firmcontact}{}
\newcommand{\firmofficeone}{}    \newcommand{\firmofficetwo}{}

\newcommand{\applyletterhead}{%
  \pagestyle{fancy}%
  \fancyhf{}%
  \renewcommand{\headrulewidth}{0pt}%
  \renewcommand{\footrulewidth}{0pt}%
  % Two partner blocks stacked, each four lines, plus the rule between them.
  % Measured rather than guessed: too small and the second partner's email
  % prints over the body of the notice.
  \setlength{\headheight}{108pt}%
  \setlength{\headsep}{14pt}%
  \setlength{\footskip}{56pt}%
  % The page must make room for the band the header now occupies. geometry is
  % re-run rather than set in persist-base, because a template using the base
  % without the letterhead should keep the tighter margins.
  \newgeometry{top=46mm, bottom=32mm, left=22mm, right=22mm,
               headheight=108pt, headsep=14pt, footskip=56pt}%
  %
  \fancyhead[L]{%
    \raisebox{-0.45\height}{\includegraphics[width=52mm]{\firmlogo}}%
  }%
  \fancyhead[R]{%
    \begin{minipage}[t]{62mm}
      \partnerblock{\partneronename}{\partneronerole}{\partneronephone}{\partneroneemail}\\[2pt]
      % The rule separates the two partners without a heavy divider.
      {\color{textSecondary}\rule{\linewidth}{0.4pt}}\\[3pt]
      \partnerblock{\partnertwoname}{\partnertworole}{\partnertwophone}{\partnertwoemail}
    \end{minipage}%
  }%
  \fancyfoot[C]{%
    \begin{minipage}{\textwidth}
      \centering
      {\sffamily\fontsize{7.5}{10}\selectfont
        Website : \textbf{\firmwebsite}\\
        Contact Us: \textbf{\firmcontact}\\
        Regd. Office: \textbf{\firmofficeone}\\
        \textbf{\firmofficetwo}}\\[6pt]
      {\fontsize{9}{11}\selectfont Page \thepage\ of \pageref{LastPage}}
    \end{minipage}%
  }%
}

% ── Correspondence idioms ────────────────────────────────────────────────────
% Centred, underlined, bold — the firm's house style for the notice type, the
% mode of service and the without-prejudice marker.
\newcommand{\noticebanner}[1]{%
  \begin{center}\underline{\textbf{#1}}\end{center}
}

% "14\textsuperscript{th} July 2026", right-aligned. Indian correspondence
% carries the ordinal; a bare 14/07/2026 reads as a form, not a letter.
\newcommand{\noticedate}[1]{%
  \begin{flushright}\textbf{Date:} #1\end{flushright}
}

% The subject line: bold, justified, and set in caps by the caller — a subject
% is the one part of a notice a recipient always reads.
\newcommand{\noticesubject}[1]{%
  \noindent\textbf{SUBJECT: #1}\par\vspace{2mm}
}

% A numbered section heading, bold, with the number hanging in the margin.
\newcommand{\noticesection}[2]{%
  \vspace{3mm}
  \noindent
  \makebox[0.6cm][l]{\textbf{#1.}}%
  \begin{minipage}[t]{\dimexpr\linewidth-0.6cm\relax}
    \textbf{#2}
  \end{minipage}\par
  \vspace{2mm}
}

% Body paragraphs inside a numbered section sit indented under the heading.
\newenvironment{noticebody}%
  {\begin{adjustwidth}{0.6cm}{0cm}\justifying}%
  {\end{adjustwidth}}
